\documentclass[11pt]{article}
\usepackage[margin=1in]{geometry}
\usepackage{booktabs,amsmath,amssymb,hyperref}

\title{HorizonGuard: Forecasting Whether a Drifting Agent Is Still Fixable,\\
and Certifying the Harm of Fixing It}
\author{Muzra Fatima}
\date{Draft, November 2026}

\begin{document}
\maketitle

\begin{abstract}
Monitors for long-horizon agents predict failure, but an accurate failure predictor can make agents
worse when its alarms trigger interventions that break runs that would have succeeded. We reframe
monitoring around recoverability: per prefix, would a goal re-anchoring rescue the run or disrupt it?
Using seed-paired counterfactual forks, we obtain rescue and disruption labels per prefix, train an
effect-aware intervention policy, and certify its disruption risk with run-level conformal risk control.
On DriftBench-CF, the effect-aware policy reaches a net benefit of 0.076 on fresh forks, against 0.001 for a
probability-only alarm. In a fork-controlled closed loop on an unseen environment and agent family, it cuts
violations from 60\% to 10\%, while a probability alarm reaches 33\%. Disruption stays below the
certified level in every one of 200 resampled calibration splits. We also report negative results: a
multi-state hazard model does not beat a single-event classifier on lead time, intervention response
transfers better than expected, and lead time follows no simple scaling law.
\end{abstract}

\section{Introduction}
% Motivation: the Intervention Paradox; why recoverability, not risk alone.

\section{Related work}
% Gap matrix: SafetyDrift, goal-drift probes, web-agent monitors, E-valuator, online safety monitoring,
% RiskShield, CIR, Drift No More.

\section{Problem formulation}
% States, potential outcomes, rescue and disruption, recoverability horizon, certified policy.

\section{Method}
\subsection{Fork labels}
\subsection{Multi-state hazard model}
\subsection{Effect model and policy}
\subsection{Disruption-controlled guarantees}

\section{Experiments}
% Tables are generated by `make eval` and copied to paper/tables by `make paper`.
\subsection{Effects and policy value}
\subsection{Guarantees: validity, coverage and sample complexity}
\subsection{Transfer and cold start}
\subsection{Closed-loop A/B}
\subsection{Limits: scaling and stealth drift}

\section{Negative results}

\section{Limitations and ethics}

\end{document}
